\textit{Night-observation backup variant.}
You have 30 minutes in total. The telescope and finder tube are already aligned.
Do not change the alignment. Write answers on the separate answer sheet.
Answers and calculations on the question paper are not evaluated.

\begin{description}
\item[(a)] The map depicts Aquila and marks the apparent magnitude of some
of its stars. Estimate the apparent magnitude of the star marked with
an \(\times\) (\(\theta\) Aql) to the same precision. Magnitudes on the map omit
the decimal point to avoid confusing it with a star, so 08 means 0.8 mag and
33 means 3.3 mag. Include the decimal point in your answer.
\hfill\mbox{[10 points]}
\ioaafig{2023_OBS-NIGHT_Stellar-f1.pdf}

\item[(b)] Measure the right ascension and declination of \(\theta\) Aql. You may
give right ascension in degrees or in hours and minutes and use the telescope
for any measurements. The Moon is near the point of the March equinox.
\hfill\mbox{[10 points]}

\item[(c)] Estimate the sidereal time within 15 minutes. You may use the
telescope for any measurements. Give the examiner your answer sheet to
record the time of observation. \hfill\mbox{[10 points]}

\item[(d)] Brocchi's Cluster, also called Collinder 399, Al Sufi's Cluster or
the Coathanger, is an asterism of ten stars. Six appear in a row across
\(1.3^\circ\). The stars are believed to be unbound. Map A shows its position
relative to other constellations and Map B shows it in a field of view of
about \(2^\circ\).
\begin{description}
\item[(i)] Centre Brocchi's Cluster in the telescope's field of view and ask
the examiner to confirm targeting and record it on your answer sheet.
If the examiner must target it for you, you lose the targeting points.
The finder scope may be easier to use. \hfill\mbox{[10 points]}
\item[(ii)] On Map B, draw two arrows pointing north and west on the celestial
sphere. \hfill\mbox{[5 points]}
\end{description}
\end{description}
\ioaafig{2023_OBS-NIGHT_Stellar-f2.pdf}
